\chapter*{Objects and Notation}
\addcontentsline{toc}{chapter}{Objects and Notation}

The notation in this book is easiest to read as a sequence of typed transformations.  Data do not
begin as an interaction graph, and a spectrum is not computed until an edge object has been
declared.  The main path is
\begin{center}
\fbox{\begin{minipage}{0.91\linewidth}
\centering
observations
$\longrightarrow$
estimator
$\longrightarrow$
typed edge operator
$\longrightarrow$
declared compression
$\longrightarrow$
spectrum or experiment.
\end{minipage}}
\end{center}
Every arrow can lose information.  The book repeatedly asks what survives each arrow and what new
assumption would be required to reverse its interpretation.

\section*{The recurring objects}

\begin{center}
\small
\setlength{\tabcolsep}{4pt}
\begin{tabular}{@{}>{\raggedright\arraybackslash}p{0.16\linewidth}
                >{\raggedright\arraybackslash}p{0.31\linewidth}
                >{\raggedright\arraybackslash}p{0.43\linewidth}@{}}
\toprule
Symbol & Type & Question answered \\
\midrule
$X=(X_1,\ldots,X_L)$ & sequence in $[q]^L$ & Which residue identity occurs at each position? \\
$p_i$ & reference distribution on the alphabet at site $i$ & Relative to which background are categorical contrasts centered? \\
$\mathcal H_i^0$ & $(q-1)$-dimensional centered categorical fiber & Which source and target perturbations count as genuine contrasts? \\
$P_i$ & projection onto $\mathcal H_i^0$ & Which constant and one-body directions are removed from a table? \\
$G_{i\leftarrow j}$ & reference-centered $q\times q$ kernel & Which source-category contrasts are coupled to which target-category contrasts? \\
$\mathcal K_{i\leftarrow j}$ & map $\mathcal H_j^0\to\mathcal H_i^0$ induced by $G_{i\leftarrow j}$ & How does a centered contrast at $j$ act on centered predictions at $i$? \\
$\mathbb K$ & block operator on $\bigoplus_i\mathcal H_i^0$ & How are all local maps assembled without scalarizing them? \\
$G_{i\leftarrow j}(C)$ & categorical kernel conditional on background $C$ & Does the effective relation change with the surrounding sequence? \\
$\overline G,\delta G$ & mean operator and context-dependent residual & What is the best static approximation, and what does it leave unexplained? \\
$Q_{ij}$ & pair probability law with declared marginals & Does a finite normalized pair law realize the local score geometry? \\
$A,L_G$ & scalar shadow and a declared graph Laplacian & Which positional organization remains after categorical information is compressed? \\
$\epsilon_{ij}$ & projected experimental epistasis table & Does a matched intervention support the predicted categorical relation? \\
\bottomrule
\end{tabular}
\end{center}

The same letter can appear with a superscript identifying its origin, such as
$G^{\mathrm P}$ for a Potts estimate or $G^{\mathrm{pLM}}$ for a language-model response.  An arrow
$j\to i$ uses $j$ as source and $i$ as target.  Static pair energies may eventually identify the
two directions, but directed responses are kept separate until reciprocity has been tested.

The distinction between $G$ and $\mathcal K$ is representational, not a second scientific claim.
In the declared categorical bases, $G_{i\leftarrow j}$ is the matrix kernel of
$\mathcal K_{i\leftarrow j}$.  The calligraphic symbol emphasizes the typed input and output;
the plain capital emphasizes the entries that are estimated, projected, plotted, or compressed.

\section*{How the notation carries type}

The typography follows a small grammar.  It does not make every symbol globally unique, but it
narrows the possibilities before an equation is read in detail.
\begin{center}
\small
\begin{tabular}{@{}>{\raggedright\arraybackslash}p{0.23\linewidth}
                >{\raggedright\arraybackslash}p{0.28\linewidth}
                >{\raggedright\arraybackslash}p{0.39\linewidth}@{}}
\toprule
Form & Usual role & Examples \\
\midrule
Lower-case italic & indices, states, scalar functions, or vectors & $i,j,a,b,t$; $f,u,v$ \\
Plain capital & finite table, matrix, random object, or model state & $T,J,G,A,P,X$ \\
Calligraphic capital & space, set, typed map, or functional & $\mathcal H,\mathcal U,\mathcal K,\mathcal E$ \\
Blackboard/fraktur capital & assembled system or expectation operator & $\mathbb K,\mathfrak G,\mathbb E$ \\
Arrow and indices & source, target, and level of action & $G_{i\leftarrow j}$ maps source $j$ toward target $i$ \\
\bottomrule
\end{tabular}
\end{center}

Decorations have stable meanings throughout the main construction:
\begin{equation*}
\begin{aligned}
\widehat G&:\ \text{sample estimate}, &
\overline G&:\ \text{declared background average},\\
\delta G&:\ G-\overline G, &
(\cdot)^*&:\ \text{population target or adjoint, as stated locally}.
\end{aligned}
\end{equation*}
Superscripts such as ${\mathrm P}$, ${\mathrm{pLM}}$, ${\mathrm{stat}}$, and ${\mathrm{dyn}}$
name origin or role; they do not change the domain and codomain unless the text explicitly says so.

Four overloaded letters remain because they are standard in the literatures being compared.  Their
syntax disambiguates them.  $L$ is sequence length in $[q]^L$ and an index bound, whereas $L_G$ or a
subscripted $L_{\mathrm{sym}},L_{\mathrm{rw}}$ is a graph Laplacian.  $P_i$ is categorical
centering, bare $P$ in the routing chapters is a row-stochastic matrix, and $P(\cdot)$ denotes a
probability.  $Q_{ij}$ is a pair law, while $Q_i$ is the Euclidean form of a categorical projector;
nuisance residualizers are written $R$ rather than adding another meaning to $Q$.  Finally, bare
$G$ can denote a generic graph in a quoted graph-network formalism, while an indexed
$G_{i\leftarrow j}$ is always a categorical kernel in this book.

\clearpage
\section*{Three spaces that should not be merged}

The book moves among three geometries:
\begin{center}
\small
\begin{tabular}{@{}>{\raggedright\arraybackslash}p{0.22\linewidth}
                >{\raggedright\arraybackslash}p{0.32\linewidth}
                >{\raggedright\arraybackslash}p{0.36\linewidth}@{}}
\toprule
Geometry & Native object & Typical operation \\
\midrule
Observation space & MSA rows, positions, and residue categories & weighting, residualization, covariance estimation \\
Categorical operator space & centered contrasts and maps between local fibers & gauge projection, reference transport, response audit \\
Position-graph space & one scalar or transported vector state per residue position & Laplacian spectrum, heat filtering, wavelets \\
\bottomrule
\end{tabular}
\end{center}
An eigenvector across MSA rows, a categorical constant mode, and an eigenvector across residue
positions are therefore different objects even when all three are called modes.

\section*{A first and a second reading}

The main argument can be followed from the opening prose of each chapter, the boxed equations,
the worked cases, and the codas.  The theorem and proof blocks establish exactness but need not be
read on the first pass.  Long derivations labeled as mathematical interludes answer a narrower
question and can likewise be deferred without losing the book's conceptual sequence.

On a second reading, each proposed correspondence can be checked with the same ledger:
\begin{equation*}
\begin{gathered}
\text{sample space},\quad
\text{normalization axis},\quad
\text{available context},\\
\text{gauge or symmetry},\quad
\text{state dynamics},\quad
\text{estimator},\quad
\text{validation endpoint}.
\end{gathered}
\end{equation*}
These are not seven hurdles that every useful model must clear.  They are seven coordinates that
prevent a valid statement in one space from being silently promoted into a stronger statement in
another.
